\section*{The Appointed Formulary in Full}

This sheet prints the complete proper of \latin{DOMINICA NONA post
Pentecosten}, \latin{II classis}, from the 1962 Vatican typical Missal,
printed pp.~388--389. Latin was collated in the study leaf's
\nolinkurl{propers/verified.md}. Scriptural English is the public-domain
Douay--Rheims revised by Challoner; the orations use the anonymous
public-domain 1861 Cummiskey hand missal. Neither English witness is an
official translation of the 1962 book. A mismatch is identified rather than
silently repaired.

\begingroup\small
\fulltextelement{1. Introit}{\latin{Antiphona ad Introitum} --- Ps. 53, 6--7; 3}
Ecce Deus ádiuvat me, et Dóminus suscéptor est ánimæ meæ: avérte mala
inimícis meis, et in veritáte tua dispérde illos, protéctor meus, Dómine.
\textnormal{Ps. ibid., 3} Deus, in nómine tuo salvum me fac: et in virtúte tua
líbera me. \textnormal{℣.~Glória Patri.}
\finishfulltextelement
\begin{fulltextenglish}{Douay--Rheims (Challoner), Ps. 53:6--7 and 53:3}
For behold God is my helper: and the Lord is the protector of my soul. Turn
back the evils upon my enemies; and cut them off in thy truth. \textit{Ps.}
Save me, O God, by thy name, and judge me in thy strength.
\end{fulltextenglish}
\fulltextgap{The missal adds \latin{protéctor meus, Dómine}; sings
\latin{líbera me} where the Clementine and Douay have ``judge me''; and omits
the psalm's opening causal particle. The doxology belongs to the Ordinary and
is not rendered here.}

\fulltextelement{2. Collect}{\latin{Oratio}}
Páteant aures misericórdiæ tuæ, Dómine, précibus supplicántium: et, ut
peténtibus desideráta concédas; fac eos, quæ tibi sunt plácita, postuláre.
Per Dóminum.
\finishfulltextelement
\begin{fulltextenglish}{Cummiskey hand missal (1861), Collect, p.~414}
May the ears of thy mercy, O Lord, be open to the prayers of thy suppliants:
and, that they may succeed in their desires, make them ask those things that
are agreeable to thee. Thro'.
\end{fulltextenglish}
\fulltextgap{The Latin asks that God grant desired things to petitioners; the
English shifts the subject and says that petitioners ``succeed.'' The Latin
therefore controls the grammatical argument in the commentary.}

\fulltextelement{3. Epistle}{\latin{Léctio Epístolæ beáti Pauli Apóstoli ad Corínthios} --- 1 Cor. 10, 6--13}
Fratres: Non simus concupiscéntes malórum, sicut et illi concupiérunt. Neque
idolólatræ efficiámini, sicut quidam ex ipsis: quemádmodum scriptum est:
Sedit pópulus manducáre et bíbere, et surrexérunt lúdere. Neque fornicémur,
sicut quidam ex ipsis fornicáti sunt, et cecidérunt una die vigínti tria
mília. Neque tentémus Christum, sicut quidam eórum tentavérunt, et a
serpéntibus periérunt. Neque murmuravéritis, sicut quidam eórum
murmuravérunt, et periérunt ab exterminatóre. Hæc autem ómnia in figúra
contingébant illis: scripta sunt autem ad correptiónem nostram, in quos fines
sæculórum devenérunt. Itaque qui se exístimat stare, vídeat ne cadat.
Tentátio vos non apprehéndat, nisi humána: fidélis autem Deus est, qui non
patiétur vos tentári supra id quod potéstis, sed fáciet étiam cum tentatióne
provéntum, ut possítis sustinére.
\finishfulltextelement
\begin{fulltextenglish}{Douay--Rheims (Challoner), 1 Cor. 10:6--13}
\dots\ that we should not covet evil things, as they also coveted. Neither
become ye idolaters, as some of them, as it is written: The people sat down to
eat and drink and rose up to play. Neither let us commit fornication, as some
of them committed fornication: and there fell in one day three and twenty
thousand. Neither let us tempt Christ, as some of them tempted and perished by
the serpents. Neither do you murmur, as some of them murmured and were
destroyed by the destroyer. Now all these things happened to them in figure:
and they are written for our correction, upon whom the ends of the world are
come. Wherefore, he that thinketh himself to stand, let him take heed lest he
fall. Let no temptation take hold on you, but such as is human. And God is
faithful, who will not suffer you to be tempted above that which you are able:
but will make also with temptation issue, that you may be able to bear it.
\end{fulltextenglish}
\fulltextgap{The missal replaces the opening words of verse 6 with the
liturgical address \latin{Fratres}; the ellipsis marks the displaced Douay
words rather than translating the missal's addition.}

\fulltextelement{4. Gradual}{\latin{Graduale} --- Ps. 8, 2}
Dómine Dóminus noster, quam admirábile est nomen tuum in univérsa terra!
\textnormal{℣.} Quóniam eleváta est magnificéntia tua super cælos.
\finishfulltextelement
\begin{fulltextenglish}{Douay--Rheims (Challoner), Ps. 8:2}
O Lord, our Lord, how admirable is thy name in the whole earth! ℣.~For thy
magnificence is elevated above the heavens.
\end{fulltextenglish}

\fulltextelement{5. Alleluia}{Ps. 58, 2}
Allelúia, allelúia. \textnormal{℣.} Eripe me de inimícis meis, Deus meus: et
ab insurgéntibus in me líbera me. Allelúia.
\finishfulltextelement
\begin{fulltextenglish}{Douay--Rheims (Challoner), Ps. 58:2}
Alleluia, alleluia. ℣.~Deliver me from my enemies, O my God; and defend me
from them that rise up against me. Alleluia.
\end{fulltextenglish}

\fulltextelement{6. Gospel}{\latin{✠ Sequéntia sancti Evangélii secúndum Lucam} --- Luc. 19, 41--47}
In illo témpore: Cum appropinquáret Iesus Ierúsalem, videns civitátem, flevit
super illam, dicens: Quia si cognovísses et tu, et quidem in hac die tua, quæ
ad pacem tibi, nunc autem abscóndita sunt ab óculis tuis. Quia vénient dies
in te: et circúmdabunt te inimíci tui vallo, et circúmdabunt te: et
coangustábunt te úndique: et ad terram prostérnent te, et fílios tuos, qui in
te sunt, et non relínquent in te lápidem super lápidem: eo quod non
cognóveris tempus visitatiónis tuæ. Et ingréssus in templum, cœpit eícere
vendéntes in illo, et eméntes, dicens illis: Scriptum est: Quia domus mea
domus oratiónis est. Vos autem fecístis illam spelúncam latrónum. Et erat
docens cotídie in templo.
\finishfulltextelement
\begin{fulltextenglish}{Douay--Rheims (Challoner), Luke 19:41--47a}
And when he drew near, seeing the city, he wept over it, saying: If thou also
hadst known, and that in this thy day, the things that are to thy peace: but
now they are hidden from thy eyes. For the days shall come upon thee: and thy
enemies shall cast a trench about thee and compass thee round and straiten
thee on every side, And beat thee flat to the ground, and thy children who
are in thee. And they shall not leave in thee a stone upon a stone: because
thou hast not known the time of thy visitation. And entering into the temple,
he began to cast out them that sold therein and them that bought. Saying to
them: It is written: My house is the house of prayer. But you have made it a
den of thieves. And he was teaching daily in the temple.
\end{fulltextenglish}
\fulltextgap{The missal supplies \latin{In illo témpore} and names Jesus and
Jerusalem where the biblical verse begins simply ``And when he drew near.''
It also ends midway through verse 47 before the opponents' response. The
printed rubric following the lesson is \latin{Credo}; the Creed belongs to the
Ordinary.}

\fulltextelement{7. Offertory}{\latin{Antiphona ad Offertorium} --- Ps. 18, 9--12}
Iustítiæ Dómini rectæ, lætificántes corda, et iudícia eius dulcióra super mel
et favum: nam et servus tuus custódit ea.
\finishfulltextelement
\begin{fulltextenglish}{Douay--Rheims (Challoner), Ps. 18:9--12; complete verses, selected fragments emphasized}
\textbf{The justices of the Lord are right, rejoicing hearts:} the commandment
of the Lord is lightsome, enlightening the eyes. The fear of the Lord is
holy, enduring for ever and ever: \textbf{the judgments of the Lord} are true,
justified in themselves. More to be desired than gold and many precious
stones: and \textbf{sweeter than honey and the honeycomb.} \textbf{For thy
servant keepeth them,} and in keeping them there is a great reward.
\end{fulltextenglish}
\fulltextgap{The Latin antiphon is a cento assembled from four fragments, not
one continuous biblical sentence. Complete English verses are shown so the
project does not manufacture its own stitched translation.}

\fulltextelement{8. Secret}{\latin{Secreta}}
Concéde nobis, quæsumus, Dómine, hæc digne frequentáre mystéria: quia,
quóties huius hóstiæ commemorátio celebrátur, opus nostræ redemptiónis
exercétur. Per Dóminum.
\finishfulltextelement
\begin{fulltextenglish}{Cummiskey hand missal (1861), Secret, p.~415}
Grant us, O Lord, we beseech thee, frequently and worthily to celebrate these
mysteries: for as many times as this commemorative sacrifice is celebrated,
so often is the work of our redemption performed. Thro'.
\end{fulltextenglish}
\fulltextgap{Both witnesses abbreviate the conclusion. The following rubric
appoints the Preface of the Most Holy Trinity, whose text belongs to the
Ordinary.}

\fulltextelement{9. Communion}{\latin{Antiphona ad Communionem} --- Io. 6, 57}
Qui mandúcat meam carnem, et bibit meum sánguinem, in me manet, et ego in eo,
dicit Dóminus.
\finishfulltextelement
\begin{fulltextenglish}{Douay--Rheims (Challoner), John 6:57}
He that eateth my flesh and drinketh my blood abideth in me: and I in him.
\end{fulltextenglish}
\fulltextgap{\latin{Dicit Dóminus} is the liturgy's attribution, not part of
the biblical verse. The missal and Douay use the Clementine number 57; many
current Bibles number the same sentence 56.}

\fulltextelement{10. Postcommunion}{\latin{Postcommunio}}
Tui nobis, quæsumus, Dómine, commúnio sacraménti, et purificatiónem cónferat,
et tríbuat unitátem. Per Dóminum.
\finishfulltextelement
\begin{fulltextenglish}{Cummiskey hand missal (1861), Postcommunion, p.~415}
May the participation of this thy sacrament, O Lord, we beseech thee, both
purify us, and unite us. Thro'.
\end{fulltextenglish}
\endgroup
